\section[Cote 86 : les cinq trièdres de l'icosaèdre et l'icosaèdre gauche]{Cote 86 : les cinq trièdres de l'icosaèdre et l'icosaèdre gauche\protect\verifie{Folder86}}\label{sec:86}

La cote 86, \emph{[Polyèdres réguliers] : notes manuscrites (s.d.)}, trente-deux pages, reprend l'icosaèdre comme un objet purement combinatoire. Sous la chemise \q{Cours}, la page~18 définit une permutation $\rho$ des quinze directions d'arêtes par le côté du pentagone d'un sommet opposé à l'arête, et la page~19 affirme $\rho^3 = \mathrm{id}$ (\lot{86}{1}{18}, \lot{86}{1}{19}). Ses orbites sont les cinq \q{trièdres distingués}. La page~20 montre que deux arêtes orthogonales n'ont pas de sommet commun et que les cinq arêtes d'un sommet ont des trièdres distincts (\lot{86}{1}{20}). Les pages~21 et~22 en tirent que le groupe $G$ des automorphismes opère sur les trièdres par permutations paires et munissent l'ensemble $T$ des trièdres d'une orientation canonique (\lot{86}{2}{21}, \lot{86}{2}{22}). Les pages~23 à~25 concluent que $G \to \mathfrak S_T$ a pour image $\operatorname{Alt}_T$ et pour noyau le sous-groupe engendré par l'antipodisme (\lot{86}{2}{23}, \lot{86}{2}{25}). Les pages~26 à~30 reconstruisent l'icosaèdre gauche $\bar\Pi = \Pi/\iota$, aujourd'hui l'hémi-icosaèdre, à partir de l'ensemble orienté $T$, et les pages~31 et~32, d'une autre encre, donnent le tableau des ensembles de drapeaux et les trois involutions sur les numérotations $(a_1, \ldots, a_5)$ (\lot{86}{2}{26}, \lot{86}{2}{30}, \lot{86}{2}{31}, \lot{86}{2}{32}).

\noindent\emph{Conventions.} Les sommets sont $\{0, \ldots, 11\}$ : $0$ et $11$ sont deux sommets opposés, $1, \ldots, 5$ et $6, \ldots, 10$ les deux pentagones de leurs voisins ; la liste des voisins est donnée une fois pour toutes, et tout le reste est calculé à partir de l'adjacence. Pour une face $\{x, y, z\}$, $\operatorname{o}(x, y, z)$ est le troisième sommet de l'autre face qui borde $\{x, y\}$ ; l'antipode $\iota x$ est l'unique sommet qui n'est ni $x$, ni un voisin de $x$, ni un voisin d'un voisin de $x$. $G$ est le groupe des bijections de l'ensemble des sommets qui conservent l'adjacence dans les deux sens. Un drapeau est un triangle orienté $(s, t, u)$ : le sommet $s$, l'arête $\{s, t\}$, la face $\{s, t, u\}$ ; il y en a $120$. Une arête est un couple de sommets adjacents, considéré à l'ordre près ; deux arêtes ont même direction si elles sont égales ou antipodales. Pour une arête $\{s, t\}$ de faces $\{s, t, u\}$ et $\{s, t, v\}$, on pose $\rho_0(s, t) = \{\operatorname{o}(s, u, t), \operatorname{o}(s, v, t)\}$, le côté $\{\Pi_{s,u}, \Pi_{s,v}\}$ de la page~18. Les numérotations de $T$ sont les bijections $\{1, \ldots, 5\} \to T$, composées à droite par les permutations des places : $(a \circ \tau)(k) = a(\tau(k))$.

\begin{enonce}[Pages 18 à 20 et 23, les trièdres distingués\piste{86-combinatorial-five-triedres}]
\begin{enumerate}
\item $\rho_0(s, t)$ est une arête et ne dépend pas de celle des deux faces que l'on appelle $u$ ; $\rho_0(t, s)$ est l'antipode de $\rho_0(s, t)$ ; si deux arêtes ont même direction, leurs images aussi. On obtient une permutation $\rho$ des quinze directions, avec $\rho^3 = \mathrm{id}$, sans point fixe, et $\rho^2\overline{\{s, t\}} = \overline{\{u, \iota v\}}$.
\item Il y a cinq orbites, les trièdres ; deux arêtes de même trièdre et de directions différentes n'ont aucun sommet commun ; en chaque sommet, les cinq arêtes ont des trièdres deux à deux distincts, et chaque trièdre est représenté. Chaque trièdre est donc une partition des douze sommets en six arêtes.
\item (Page~23) Pour une face $\{u, v, w\}$, avec $u_1 = \operatorname{o}(v, w, u)$, $v_1 = \operatorname{o}(w, u, v)$, $w_1 = \operatorname{o}(u, v, w)$, les directions $\overline{\{u, v_1\}}$, $\overline{\{v, w_1\}}$, $\overline{\{w, u_1\}}$ sont dans un même trièdre, $\overline{\{u, w_1\}}$, $\overline{\{v, u_1\}}$, $\overline{\{w, v_1\}}$ dans un autre, et ces deux trièdres et ceux des trois côtés de la face sont les cinq trièdres. Les arêtes $\{u, v\}$ et $\{v, w_1\}$ ne sont jamais dans le même trièdre.
\end{enumerate}
\end{enonce}

\begin{proof}
Ce sont des énoncés finis sur l'icosaèdre, vérifiés par le noyau de Lean (\texttt{decide}), sans évaluation compilée. Pour~(2), on étiquette les trente arêtes par $\{0, \ldots, 4\}$ et l'on vérifie que deux arêtes ont la même étiquette si et seulement si la direction de l'une est celle de l'autre, de son image par $\rho$ ou de son image par $\rho^2$ : les étiquettes sont alors les orbites, et l'étiquetage n'est qu'un nom donné à chacune.
\end{proof}

\begin{encadre}
\begin{theoreme}[Grothendieck, pages 20 à 25, l'action sur les trièdres\piste{86-combinatorial-five-triedres}]
\begin{enumerate}
\item Un automorphisme est déterminé par l'image du drapeau $(0, 1, 2)$, et tout drapeau est l'image de $(0, 1, 2)$ par un automorphisme : $G$ a $120$ éléments.
\item Tout automorphisme $g$ commute à $\iota$ et à $\rho$, et l'application qui envoie le trièdre d'une arête $a$ sur celui de $g(a)$ définit un homomorphisme $G \to \mathfrak S_T$.
\item Son image est contenue dans $\operatorname{Alt}_T$ et lui est égale.
\item Son noyau est $\{1, \iota\}$.
\end{enumerate}
\end{theoreme}
\end{encadre}

\begin{proof}
(1) Si $g$ conserve l'adjacence, $g(\operatorname{o}(x, y, z)) = \operatorname{o}(gx, gy, gz)$, car $\operatorname{o}(x, y, z)$ est l'unique voisin commun de $x$ et $y$ distinct de $z$. Les douze sommets s'obtiennent à partir de $0$, $1$, $2$ par neuf applications de $\operatorname{o}$ ; $g$ est donc déterminé par $(g0, g1, g2)$. Réciproquement, pour chaque drapeau $(x, y, z)$, la même suite d'applications de $\operatorname{o}$ définit une application dont on vérifie qu'elle est bijective et conserve l'adjacence. (2) $\iota x$ est caractérisé par l'adjacence, donc $g\iota = \iota g$. Comme $\rho_0(s, t)$ est formé de deux valeurs de $\operatorname{o}$ et ne dépend pas du choix de $u$, $g$ envoie $\rho_0(s, t)$ sur $\rho_0(gs, gt)$ ; $g$ conserve donc la relation \q{même trièdre}, et l'action sur les cinq étiquettes est bien définie et multiplicative. (3) et (4) Par~(1), tout $g$ est l'automorphisme associé à un drapeau, et la permutation de $T$ qu'il induit se calcule à partir de ce drapeau. On vérifie sur les $120$ drapeaux que ces permutations sont paires, qu'elles sont au nombre de $60$, ce qui est l'ordre de $\operatorname{Alt}_T$, et que seuls les drapeaux $(0, 1, 2)$ et $(\iota 0, \iota 1, \iota 2)$ donnent l'identité.
\end{proof}

\begin{remarque}
La démonstration de la page suit un autre chemin : la parité par les générateurs $\bar\sigma_0$, $\bar\sigma_1$, $\bar\sigma_2$ et le cas $n \equiv 1 \pmod 4$ du polygone, l'image par les ordres $6$, $10$ et $4$ des images des stabilisateurs d'une face, d'un sommet et d'une arête. Ici l'énoncé est établi par le décompte des drapeaux ; seul le fait que les symétries d'un pentagone sont des permutations paires de ses sommets est vérifié pour lui-même (énoncé suivant). La vérification confirme le calcul de $\rho^2$ et de $\rho^3$, que la page ne détaille pas, et la correction de la lecture à la page~23 : avec $\{u, v\}$, la page met dans un même triple deux arêtes de trièdres différents. Les deux choix de $u$ dans la définition de $\rho_0$ donnent la même arête, et le passage de $s$ à $t$ donne l'arête antipodale : $\rho$ est bien défini sur les directions, comme la page l'affirme.
\end{remarque}

\begin{encadre}
\begin{theoreme}[D'après les pages 21 et 26 à 32, forme complétée par l'édition, l'icosaèdre gauche par l'ensemble orienté $T$\piste{86-hemi-icosahedron-five-set-model}]
À un drapeau $(s, t, u)$, avec $v = \operatorname{o}(s, t, u)$, associons la numérotation $\nu(s, t, u)$ de $T$ formée des trièdres des arêtes $\{s, t\}$, $\{s, v\}$, $\{s, \Pi_{s,v}\}$, $\{s, \Pi_{s,u}\}$, $\{s, u\}$ : les cinq arêtes de $s$ dans l'ordre du pentagone, à partir de $t$, dans le sens qui s'éloigne de $u$ (page~30).
\begin{enumerate}
\item Pour $g \in G$, $\nu(gs, gt, gu)$ est $\nu(s, t, u)$ composée à gauche par la permutation de $T$ induite par $g$.
\item Toutes les numérotations $\nu(s, t, u)$ ont la même parité ; elles définissent donc une orientation $\varpi$ de $T$ qui ne dépend d'aucun choix. Toute numérotation compatible avec $\varpi$ est de la forme $\nu(s, t, u)$, et $\nu(s, t, u) = \nu(s', t', u')$ si et seulement si $(s', t', u')$ est $(s, t, u)$ ou $(\iota s, \iota t, \iota u)$. Les $60$ drapeaux de $\bar\Pi$ sont ainsi les $60$ numérotations compatibles avec $\varpi$.
\item Changer le sommet, l'arête ou la face d'un drapeau, c'est-à-dire $(s, t, u) \mapsto (t, s, u)$, $(s, u, t)$, $(s, t, v)$, revient à composer $\nu$ à droite par
\[ \tau_0 = (2\,4)(3\,5), \qquad \tau_1 = (1\,5)(2\,4), \qquad \tau_2 = (2\,5)(3\,4) \]
(permutations des places) : $\sigma_0(a) = (a_1, a_4, a_5, a_2, a_3)$, $\sigma_1(a) = (a_5, a_4, a_3, a_2, a_1)$, $\sigma_2(a) = (a_1, a_5, a_4, a_3, a_2)$.
\item Le $\sigma_1$ de la page~31, $(a_2, a_1, a_5, a_4, a_3)$, est $\tau_1' = (1\,2)(3\,5) = \tau_2\tau_1\tau_2$, et $\tau_2$ commute à $\tau_0$. Le triple de la page est donc l'action des trois mêmes changements sur la numérotation $(s, t, u) \mapsto \nu(s, t, v)$.
\item Avec le triple de la page, $\sigma_1\sigma_0 : a \mapsto (a_4, a_1, a_3, a_2, a_5)$, $\sigma_2\sigma_1 : a \mapsto (a_2, a_3, a_4, a_5, a_1)$ et $\sigma_0\sigma_2 : a \mapsto (a_1, a_3, a_2, a_5, a_4)$ sont d'ordres $3$, $5$ et $2$, et $\sigma_0\sigma_1$ est d'ordre $3$.
\item Deux drapeaux ont le même sommet dans $\bar\Pi$ si et seulement si leurs numérotations tracent le même pentagone sur $T$ ; la même arête si et seulement si elles donnent la même partition $\{a_1\} \sqcup \{a_2, a_3\} \sqcup \{a_4, a_5\}$ de type $(2, 2, 1)$ ; la même face si et seulement si elles donnent la même partie $\{a_1, a_3, a_5\}$.
\item Les symétries du pentagone $1\,2\,3\,4\,5$ sont des permutations paires des places. Par suite une numérotation de $T$ trace le pentagone d'un sommet de $\bar\Pi$ si et seulement si elle est compatible avec $\varpi$. On obtient $6$ pentagones, $15$ partitions de type $(2, 2, 1)$, c'est-à-dire toutes, et les $10$ parties à trois éléments.
\end{enumerate}
\end{theoreme}
\end{encadre}

\begin{proof}
(1) Les cinq arêtes de $s$ s'écrivent avec $\operatorname{o}$, que $g$ respecte, et $g$ opère sur les trièdres par l'énoncé précédent. (2) Par~(1), $\nu(g0, g1, g2)$ est la permutation de $T$ induite par $g$, composée avec $\nu(0, 1, 2)$ ; avec l'étiquetage choisi, $\nu(0, 1, 2)$ est la permutation des places $\tau_2 = (2\,5)(3\,4)$, qui est paire, et toute permutation paire s'obtient ainsi par la surjectivité sur $\operatorname{Alt}_T$. Deux drapeaux $g(0,1,2)$ et $h(0,1,2)$ ont la même numérotation si et seulement si $g^{-1}h$ est dans le noyau $\{1, \iota\}$, c'est-à-dire $h = g$ ou $h = g\iota$, et $g\iota(0,1,2) = \iota g(0,1,2)$. (3) Les changements de drapeau sont $g \mapsto g s_i$, où $s_0$, $s_1$, $s_2$ sont les automorphismes qui envoient $(0, 1, 2)$ sur $(1, 0, 2)$, $(0, 2, 1)$, $(0, 1, 5)$ ; par~(2), il suffit de vérifier les trois égalités de permutations $\pi(s_i)\,\nu(0,1,2) = \nu(0,1,2)\,\tau_i$, où $\pi(s_i)$ est la permutation de $T$ induite. (4) et (5) sont des calculs de permutations ; par~(3), $\nu(\sigma_2(\sigma_1 r)) = \nu(r)\tau_1\tau_2 = \nu(\sigma_2 r)\,\tau_2\tau_1\tau_2$, puisque $\tau_2^2 = 1$. (6) Par~(1), en remplaçant les deux drapeaux par leurs images sous $g^{-1}$, où $g$ envoie $(0, 1, 2)$ sur le premier, on se ramène au cas où le premier est $(0, 1, 2)$ : les trois conditions sur les drapeaux sont conservées par $g$, et les trois conditions sur les numérotations par la composition à gauche avec une permutation de $T$. Le cas du drapeau $(0, 1, 2)$ se vérifie sur les $120$ drapeaux. (7) On vérifie que les numérotations qui tracent le pentagone $1\,2\,3\,4\,5$ sont les dix éléments du groupe diédral et que ces dix permutations sont paires ; si deux numérotations $a$, $b$ tracent le même pentagone, $a^{-1}b$ trace le pentagone standard, donc $a$ et $b$ ont même parité. Les trois décomptes sont des calculs sur les $120$ drapeaux.
\end{proof}

\begin{remarque}
L'orientation $\varpi$ est celle des pages~21 et~22 : la parité de $\nu(s, t, u)$ ne dépend ni du drapeau ni du sens de parcours, et c'est le cas $n \equiv 1 \pmod 4$ du pentagone qui le permet, puisque ses dix symétries sont paires ; pour un hexagone, la symétrie qui échange les extrémités de deux côtés opposés est un produit de trois transpositions, donc impaire. La vérification confirme ce que la lecture signale aux pages~31 et~32. Les pages~26 à~32 donnent les bijections et les ensembles de la page~31 sans les incidences ; les incidences du point~6 et les points~4 et~5 sont de l'édition. Le $\sigma_1$ de la page n'est pas celui que donne le dictionnaire de la page~30 ; il s'en déduit par conjugaison par $\sigma_2$, et correspond au dictionnaire qui lit le pentagone de $s$ à partir de $t$ en direction de $u$. Les deux triples définissent donc des ensembles à opérateurs isomorphes. L'étiquette \q{$\sigma = \sigma_0\sigma_1$} de la page~32 accompagne la valeur de $\sigma_0\sigma_2$, d'ordre~$2$. L'étiquetage des trièdres a été choisi de sorte que $\varpi$ soit la classe des numérotations paires ; un autre étiquetage changerait la classe, non l'énoncé.
\end{remarque}

Ne sont pas démontrés ici : les axiomes des polyèdres combinatoires de la page~2 et le théorème de pleine fidélité des pages~10 à~16, l'icosaèdre étant donné par sa liste de voisins ; l'identification des trièdres aux repères orthogonaux et aux cinq octaèdres inscrits de l'icosaèdre régulier ; la décomposition $G = \{1, \iota\} \times G^{+}$ et l'isomorphisme $G^{+} \simeq \operatorname{Alt}_T$ du corollaire de la page~26, faute d'une définition de l'orientation de $\Pi$ ; la description des arêtes par les triplets $(t, \gamma, \alpha)$ de la page~27 et les incidences a) à d) des pages~28 et~29 sous cette forme, que l'on a remplacées par les partitions et les parties à trois éléments de la page~31 ; les lignes $\vec A$, $A^{\uparrow}$, $A^{\wedge}$, $\tilde S$, $\tilde A$, $\tilde F$ du tableau de la page~31 ; le fait que $\tau_0$, $\tau_1$, $\tau_2$ engendrent $\operatorname{Alt}_5$ et opèrent simplement transitivement sur les $60$ numérotations, qui suivrait de la connexité du graphe des drapeaux, non vérifiée.
